%% ma: L10-B14-0002
%% lop: 10
%% chuong: 5
%% bai: 14
%% dang: 10.14.2
%% loai: TL
%% muc: TH
%% dap_an: "-"
%% nguon: "chuyen-de-cac-so-dac-trung-cua-mau-so-lieu-khong-ghep-nhom-toan-10-chuong-trinh-moi.pdf (GV. Phan Nhật Linh) (P04), Bài 03 Dạng 1 Bài tập tự luận 2"
%% kien_thuc: "Lớp 10 Bài 13 (tứ phân vị), Bài 14 (khoảng tứ phân vị)"
%% trang_thai: chinh-thuc
%% ghi_chu: "Tự duyệt theo yêu cầu giáo viên 22/10/2026."
\begin{ex}
Hãy tính khoảng biến thiên và khoảng tứ phân vị của mẫu số liệu: $15$; $20$; $1$; $2$; $4$; $3$; $7$; $5$.
\loigiai{Sắp xếp: $1$; $2$; $3$; $4$; $5$; $7$; $15$; $20$. $R=20-1=19$.

$Q_2=\dfrac{4+5}2=4{,}5$; $Q_1=\dfrac{2+3}2=2{,}5$; $Q_3=\dfrac{7+15}2=11$.

$\Delta_Q=11-2{,}5=8{,}5$.}
\end{ex}
